\section{A Game-Theoretic Framework for Modular Specifications}
\label{sec:framework}

\subsection*{A Spec-Agnostic Semantic Lens for Architectural Analysis}

\paragraph{Abstract.}
We present a semantic, specification-agnostic framework for analyzing modular software architectures through \emph{multi-player coordination games}.
Component implementations act as strategies; the composed system induces demonic nondeterministic behaviors; and local specifications define binary obligation-satisfaction utilities.
A chosen solution concept~$\mathsf{SC}$ selects architecturally relevant profiles.

Architectural properties such as coupling, substitutability, minimality, and information hiding arise by comparing these $\mathsf{SC}$-selected profiles with globally desirable behaviors. As a motivating illustration rather than a general theorem (see \Cref{subsec:pne-caveat} below), some solution concepts---notably Pure Nash Equilibrium (PNE)---can select equilibrium profiles that are \textbf{undesired} stability traps: points where no component can locally repair its obligation even though the global system behavior is undesirable.

The approach is semantic, lightweight, and applies to any specification formalism with a behavioral interpretation.

\subsection{Motivation}

In modular architectures, components interact through a shared semantic environment while each carries its own local obligations.
Models such as game semantics, assume--guarantee/rely-guarantee
reasoning~\cite{Jones1983RelyGuarantee,AbadiLamport1993ComposingSpecifications},
interface theories~\cite{deAlfaroHenzinger2001InterfaceAutomata}, or
contract-based design~\cite{Benveniste2018ContractsForSystemDesign} typically
use \emph{two-player} perspectives (a component against its environment).
Some approaches (MPST, certain assume-guarantee frameworks~\cite{ChiltonAssumeGuarantee}) extend to multi-role perspectives, but here we adopt an even more general view:

\begin{quote}
\emph{A tuple of specifications induces a multi-player game.
Implementations act as strategies; the composed system induces behaviors; a solution concept selects key profiles for architectural analysis.}
\end{quote}

Crucially, the framework is \emph{agnostic} about what those profiles \emph{mean}: depending on the solution concept, they may represent desirable configurations or undesirable ones.

\subsection{System Model}

Let $C=\{c_0,\dots,c_{n-1}\}$ be the set of components.

\paragraph{Specifications and Behaviors.}

\begin{itemize}
\item A behavior domain $\mathsf{Beh}$.
\item For each component $c$:
  \begin{itemize}
  \item a specification space $\mathsf{Spec}_c$,
  \item an implementation space $\Sigma_c$,
  \item a semantics $\llbracket - \rrbracket_c : \mathsf{Spec}_c \to \mathcal P(\mathsf{Beh})$.
  \end{itemize}
\end{itemize}

The satisfaction set may be empty.
If $\llbracket S_c \rrbracket_c = \varnothing$, then component $c$ is permanently unsatisfiable: $u_c(\sigma)=0$ for all~$\sigma$.

\paragraph{Demonic Composition Semantics.}

A joint profile is $\sigma=(\sigma_c)_{c\in C}$.

\[
\mathsf{out} : \prod_{c\in C}\Sigma_c \to \mathcal P(\mathsf{Beh})
\]

$\mathsf{out}$ is \emph{demonic}: every behavior in $\mathsf{out}(\sigma)$ must satisfy obligations.
Outputs are assumed non-empty.
Utilities therefore depend on the \emph{entire} behavior set.

This demonic reading is appropriate for safety-style reasoning (``every possible behavior must be safe'') and pessimistic architectural analysis. For models of intentional internal nondeterminism or optimization under uncertainty, one would instead use angelic or probabilistic readings; these can be accommodated by changing the definition of $u_c$ and the behavior model, but we do not pursue that here.

\paragraph{Observation Granularity.}

Each component observes via
\[
\pi_c : \mathsf{Beh} \to \mathsf{Obs}_c,
\qquad
\pi_c(B) = \{\pi_c(b) \mid b \in B\}.
\]

The projection $\pi_c$ may collapse behaviors; it defines the abstraction level and visible nondeterminism.

\paragraph{Global Desirability.}

A global correctness predicate:
\[
\mathsf{Valid} \subseteq \mathsf{Beh},
\qquad
\mathsf{ValidBeh} = \{\sigma \mid \mathsf{out}(\sigma) \subseteq \mathsf{Valid}\}.
\]

No relationship between $\mathsf{Valid}$ and local obligations is assumed.

\subsection{The Induced Game}

Given $S=(S_c)_{c\in C}$:

\begin{itemize}
\item players: components;
\item strategies: implementations;
\item outcomes: $\mathsf{out}(\sigma)$;
\item utilities:
\[
u_c(\sigma)=
\begin{cases}
1 & \text{if }\mathsf{out}(\sigma)\subseteq\llbracket S_c\rrbracket_c\\
0 & \text{otherwise.}
\end{cases}
\]
\end{itemize}

Utilities are binary, demonic, and obligation-based: any violating behavior forces utility~0, and satisfying behaviors are not ranked.

\subsection{Solution Concepts}

A solution concept gives a set
\[
\mathsf{SC}(S) \subseteq \prod_{c\in C} \Sigma_c
\]
of architecturally relevant profiles.

Different solution concepts correspond to different architectural questions.
Pure Nash Equilibrium (PNE) captures unilateral repairability;
dominant strategies capture robustness to any environment behavior;
max-min or worst-case concepts capture hardening against adversarial contexts;
correlated or cooperative equilibria model designed coordination.
In this paper we work explicitly with PNE as a running example, but the definitions are parametric in SC and can be instantiated with these standard notions.


\paragraph{Pure Nash Equilibrium (PNE).}

Let $(\sigma_{-c}, \sigma'_c)$ replace the $c$-component of $\sigma$.

\[
\sigma \in \mathsf{PNE}(S)
\;\Longleftrightarrow\;
\neg \exists c,\sigma'_c.\;
u_c(\sigma_{-c},\sigma'_c) > u_c(\sigma).
\]

Because utilities are binary and demonic:
\begin{itemize}
\item only $0\to 1$ improvements matter;
\item profiles where all utilities are $0$ may still be PNE;
\item PNE can identify \emph{undesired stability traps}.
\end{itemize}

This undesirability is specific to PNE under binary demonic utilities; other solution concepts may select desirable configurations.

\label{subsec:pne-caveat}
\paragraph{Caveat.} The claim that PNE ``identifies'' undesired stability traps is an informal observation motivating the framework, illustrated by the all-zero-utility example above, not a general theorem: we have not characterized which specification tuples $S$ admit an all-desirable PNE, which admit only trap-like PNE, or how common each case is. Formalizing that characterization (e.g.\ as a condition on $\mathsf{Valid}$, $\mathsf{out}$, and the $u_c$) is left open; until it exists, this subsection should be read as a worked motivating example for the framework's vocabulary rather than as evidence that PNE analysis of module specifications reliably surfaces bad states.

Before treating that characterization as a fresh open problem, it should be
checked against rational verification's existing complexity theory (see the
paragraph below), which already studies ``which games admit desirable
equilibria'' in a closely related
form~\cite{FismanKupfermanLustig2010RationalSynthesis,AbateEtAl2021RationalVerification};
``self-contained'' is precisely the property that makes a question likely to
already exist in a neighboring literature.

\paragraph{Named intermediate solution concepts.}
Forceability-style coupling (\Cref{prop:indep-is-sc-all} below) is sound and
compositional but paranoid: it certifies safety by assuming every other
component behaves adversarially, discarding any information about what the
other components actually are or want. PNE is close to the opposite failure:
it can fail to rule out profiles that are not even individually compliant.
The formal-methods and multi-agent-systems literature has already engineered
several points strictly between these extremes, each trading paranoia for
compositionality in a controlled way:
\begin{itemize}[nosep]
  \item \textbf{Secure equilibria}~\cite{ChatterjeeHenzingerJurdzinski2006SecureEquilibria}
  use lexicographic objectives (primarily achieve your own goal, secondarily
  hurt the opponent) to obtain equilibria with a unique value, engineered
  specifically to support compositional verification --- a targeted fix for
  the stability-trap problem of \Cref{subsec:pne-caveat}.
  \item \textbf{Assume-admissible synthesis}~\cite{BrenguierRaskinSankur2017AssumeAdmissible}
  has each component assume the others play \emph{admissible} (non-dominated)
  strategies, rather than assuming either full adversarial behavior or a
  fixed-point equilibrium. This composes, unlike PNE, while being strictly
  less paranoid than forceability's $\forall\sigma_{-D}$.
  \item \textbf{Rational synthesis / rational
  verification}~\cite{FismanKupfermanLustig2010RationalSynthesis,AbateEtAl2021RationalVerification}
  certify a specification under the assumption that the environment plays
  \emph{some} equilibrium of its own objectives, with a developed complexity
  theory for when this is decidable.
\end{itemize}
Ordering these by how much they assume about the environment gives a monotone
family, in the sense of SC Monotonicity below, sitting between
$\mathsf{SC}_{\mathrm{All}}$ (\Cref{prop:indep-is-sc-all}: assumes nothing
beyond raw compliance) and forceability (assumes the worst). Each of these is
a solution concept in this appendix's sense and generates its own coupling,
substitutability, and minimality notions by the schema above; we have not
carried out that instantiation, nor checked which points of this spectrum
survive \KFOL's symmetry-group action (\Cref{sec:symmetry}) or remain
two-coalition facts under coalition merging
(\Cref{subsec:coalitional-forceability}). Both are left open.

\paragraph{SC Monotonicity.}

If $\mathsf{SC}_1(S)\subseteq\mathsf{SC}_2(S)$, then
\[
\mathsf{SC}_2(S) \subseteq \mathsf{ValidBeh}
\;\Rightarrow\;
\mathsf{SC}_1(S) \subseteq \mathsf{ValidBeh}.
\]

No monotonicity holds for specification weakening or behavior refinement.

\subsubsection{Implementation Independence}

The coupling notions above are defined \emph{relative to a solution concept}:
they only quantify over profiles in $\mathsf{SC}(S)$.
Sometimes we want a stronger, \emph{solution-concept-independent} notion that
says: once a set of components meets its local obligations, the system is
globally well-behaved, regardless of which particular implementations were
chosen.

Let $D \subseteq C$ be a set of components.
We say that the specification tuple $S$ is
\emph{implementation independent on $D$} if any profile that satisfies
the obligations of all components in $D$ is globally desirable:
\[
  \mathsf{Indep}_D(S)
  \;\triangleq\;
  \forall \sigma \in \textstyle\prod_{c\in C} \Sigma_c.\;
    \biggl(
      \bigwedge_{c\in D} u_c(\sigma) = 1
    \biggr)
    \;\Longrightarrow\;
    \mathsf{out}(\sigma) \subseteq \mathsf{Valid}.
\]

Intuitively, $\mathsf{Indep}_D(S)$ states that implementations for the
components in $D$ can be chosen \emph{independently}: any combination of
implementations that meets their local specifications yields globally
desirable behavior, irrespective of the particular choices made for the
remaining components in $C \setminus D$ (provided they are included in
$\sigma$ and participate in $\mathsf{out}$ as usual).

\paragraph{Reading via assume-guarantee contracts.}
This quantifier shape --- $\forall$ over $D$'s own compliant implementations,
$\forall$ over the environment --- is exactly the satisfaction relation of
assume-guarantee/rely-guarantee
reasoning~\cite{Jones1983RelyGuarantee,AbadiLamport1993ComposingSpecifications},
interface theories~\cite{deAlfaroHenzinger2001InterfaceAutomata}, and their
modern synthesis in contract-based
design~\cite{Benveniste2018ContractsForSystemDesign}: a component's contract
is satisfied demonically on \emph{both} sides, not only the environment's.
Forceability (the main text's \Cref{def:k-force}) is angelic on the player's
own side ($\exists\profile_i$: some strategy works) and demonic on the
environment's side ($\forall\profile_{-i}$); $\mathsf{Indep}_D$ is demonic on
\emph{both} sides ($\forall\sigma_D$ compliant, $\forall\sigma_{-D}$), which
is exactly what makes it a statement about the \emph{contract itself} rather
than about the existence of a witness --- you should not need to find the
clever compliant implementation, any of them must do. That body of theory
already develops the composition operators (quotient/residuation), refinement
orders, and compatibility conditions this appendix does not; what follows
should be read as an instance of that theory inside this paper's vocabulary,
not as new theory.

Two special cases are worth highlighting:

\begin{itemize}
  \item \textbf{Pairwise independence.}
  For two components $a,b \in C$, implementation independence on
  $D = \{a,b\}$ says that whenever both $a$ and $b$ satisfy their local
  obligations in some profile, the resulting global behavior is valid.
  This matches the informal idea that the modules for $a$ and $b$ can be
  implemented and combined without introducing architectural failures.

  \item \textbf{Full independence.}
  Implementation independence on $D = C$ says that as soon as
  \emph{all} components satisfy their local specifications in a profile,
  the system is globally desirable:
  \[
    \mathsf{Indep}_C(S)
    \;\equiv\;
    \forall \sigma.\;
      \biggl(
        \bigwedge_{c\in C} u_c(\sigma) = 1
      \biggr)
      \;\Longrightarrow\;
      \mathsf{out}(\sigma) \subseteq \mathsf{Valid}.
  \]
  This is a very strong architectural property:
  any locally compliant joint implementation is guaranteed to be globally
  acceptable.
\end{itemize}

Unlike the $\mathsf{SC}$-relative coupling notion
$\mathsf{SC}(S) \subseteq \mathsf{ValidBeh}$, implementation independence
quantifies over \emph{all} profiles that satisfy the obligations of the
chosen components $D$, regardless of whether they are selected by the
solution concept.
It therefore captures a stricter, solution-concept-agnostic form of
``independent implementability''.

\paragraph{A compatibility (nonemptiness) side condition.}
As stated, $\mathsf{Indep}_D(S)$ is vacuously true whenever $D$'s
locally-compliant profile set is empty --- i.e.\ whenever $D$'s own local
specification is unsatisfiable in isolation. This is the same vacuity
contract theory calls out under the name \emph{compatibility}. We do not
build a nonemptiness requirement into the definition above, so that
$\mathsf{Indep}_D(S)$ remains a pure implication matching the rest of this
appendix's style; but any use of $\mathsf{Indep}_D(S)$ as evidence of
\emph{safe} modularity, rather than merely \emph{vacuously consistent}
modularity, should check nonemptiness of $D$'s compliance set separately.

\begin{proposition}[Full independence is coupling at the compliance solution concept]
\label{prop:indep-is-sc-all}
Define $\mathsf{SC}_{\mathrm{All}}(S) := \{\sigma \in \prod_c \Sigma_c :
\forall c\in C,\ u_c(\sigma)=1\}$, the degenerate solution concept that
imposes no strategic or rationality assumption beyond raw local compliance.
Then
\[
\mathsf{Indep}_C(S) \iff \mathsf{SC}_{\mathrm{All}}(S) \subseteq \mathsf{ValidBeh}.
\]
Moreover, if $\sigma$ is a profile in which every $c\in C$ plays a dominant
strategy for its own local specification (i.e.\ $u_c(\sigma_{-c},\sigma_c)=1$
for every $\sigma_{-c}$), then $\sigma\in\mathsf{SC}_{\mathrm{All}}(S)$; so
the dominant-strategy solution concept is a \emph{subset} of
$\mathsf{SC}_{\mathrm{All}}$, not in general equal to it (a dominant strategy
is compliant in every environment, but a profile can be compliant in the one
environment actually realized without any component's strategy being
dominant). By SC Monotonicity, $\mathsf{Indep}_C(S)$ therefore implies loose
coupling under the dominant-strategy solution concept, of which forceability
is the binary-utility instance --- but the implication runs only in that
direction.
\end{proposition}
\begin{proof}
Immediate from unfolding both definitions: $\mathsf{Indep}_C(S)$ says
$\forall\sigma.\ (\bigwedge_c u_c(\sigma)=1) \Rightarrow
\mathsf{out}(\sigma)\subseteq\mathsf{Valid}$, which is literally
$\mathsf{SC}_{\mathrm{All}}(S)\subseteq\mathsf{ValidBeh}$ once
$\mathsf{SC}_{\mathrm{All}}(S)$ is unfolded. The dominant-strategy inclusion
and the monotonicity consequence are as stated above.
\end{proof}

\begin{remark}[Why PNE escapes this hierarchy]
$\mathsf{PNE}(S)$ is \emph{not} in general a subset of
$\mathsf{SC}_{\mathrm{All}}(S)$: as already noted, an all-zero-utility profile
can be a PNE. So \Cref{prop:indep-is-sc-all}'s monotonicity argument says
nothing about PNE-coupling in either direction --- $\mathsf{Indep}_C(S)$ can
hold while some PNE is still an undesired stability trap, because that PNE
need not be locally compliant at all and so falls entirely outside what
$\mathsf{Indep}_C$ constrains. This is the precise mechanism behind
\Cref{subsec:pne-caveat}'s caveat, not merely an informal observation about
it.
\end{remark}

\begin{remark}[A tempting but likely-flat $K$-ary candidate]
Byzantine-style robustness --- ``valid under arbitrary deviation by any
$\le f$ of the $K$ components'' --- is the notion in this framework that most
looks like genuine $K$-ary content, since it quantifies over a whole family
of coalitions rather than one fixed coalition. But it is exactly a finite
conjunction of two-coalition facts, one forceability-style statement per
coalition of size $\le f$, and the flattening result of
\Cref{prop:coalition-merge-lossless} applies to each conjunct separately.
This is the third time in this paper that a plausible-looking source of
$K$-ary content collapses into a family of two-coalition facts, after the
tensor/coalition-merging discussion of \Cref{subsec:open-games} and the
IF-binder conjecture of \Cref{subsec:coalitional-forceability}. If this
calculus has genuine $K\ge 3$ content anywhere, the pattern across all three
attempts points the same way: it will come from information asymmetry
between components, not from counting coalitions.
\end{remark}

\subsection{Architectural Notions}

Architectural properties compare $\mathsf{SC}(S)$ with desirable profiles.

\paragraph{Coupling.}

A system is \emph{loosely coupled} under $\mathsf{SC}$ when
\[
\mathsf{SC}(S)\subseteq \mathsf{ValidBeh}.
\]

Otherwise some SC-profile is globally undesirable.
If $\mathsf{SC}(S)=\varnothing$, the condition holds vacuously.

\paragraph{Substitutability Relative to the Global Property.}

Fix component $B$.
For all environments $\sigma_{-B}$ and all $\sigma_B,\sigma'_B$ with
$u_B(\sigma_{-B},\sigma_B)=u_B(\sigma_{-B},\sigma'_B)=1$:

\[
\mathsf{out}(\sigma_{-B},\sigma_B)\subseteq\mathsf{Valid}
\;\Longleftrightarrow\;
\mathsf{out}(\sigma_{-B},\sigma'_B)\subseteq\mathsf{Valid}.
\]

Quantification is over \emph{all} environments, not only $\mathsf{SC}$-reachable ones.
This expresses $\mathsf{Valid}$-preserving substitutability, not refinement or observational equivalence.

\paragraph{Minimality of Specifications.}

Semantic weakening:
\[
S'_c \le S_c \Rightarrow 
\llbracket S_c\rrbracket_c \subseteq \llbracket S'_c\rrbracket_c.
\]

$S_B$ is \emph{$\mathsf{SC}$-minimal} in context $S_A$ when:
\begin{enumerate}
\item $\mathsf{SC}(S)\neq\varnothing$ and $\mathsf{SC}(S)\subseteq\mathsf{ValidBeh}$;
\item no strictly weaker $S'_B < S_B$ preserves \emph{all} desirable SC-stable profiles.
\end{enumerate}

\paragraph{Information Hiding (Stability-Oriented).}

\[
\sigma_{-c}=\sigma'_{-c}
\ \land\
\pi_c(\mathsf{out}(\sigma))=\pi_c(\mathsf{out}(\sigma'))
\;\Rightarrow\;
(\sigma\in\mathsf{SC}(S)\Rightarrow \sigma'\in\mathsf{SC}(S)).
\]

A strict extensional criterion:
equality of projected behavior sets is required; adding new projected behaviors or splitting one into multiple distinguishable observations breaks hiding.
This is stricter than trace- or simulation-based hiding unless accounted for in~$\pi_c$.
\subsection{Adversarial Test Components}
\label{subsec:adversarial-tests}

In a modular setting, \emph{tests} are naturally modeled as \emph{adversarial}
components: they are written with the explicit intention of catching bad
behavior of other components. This is slightly at odds with everyday testing
terminology: here a ``good'' test is one that is \emph{bad for} at least one
incorrect implementation.

Let $\mathcal I$ be the set of implementations of a given interface or
component role, and let $\mathcal G \subseteq \mathcal I$ be the set of
\emph{intended} implementations. For example, one may take $\mathcal G$ to be those
implementations that both satisfy a local specification $S_c$ and lead to
globally desirable behavior (with respect to $\mathsf{Valid}$) under some
chosen class of environments.

\begin{definition}[Test suite]
Let $\mathcal I$ be the implementation space for a given interface, and let
$\mathcal G \subseteq \mathcal I$ be the intended implementations. A
\emph{test suite} is a function
\[
  T : \mathcal I \to \{\mathsf{pass}, \mathsf{fail}\}.
\]
It is
\begin{itemize}
  \item \emph{non-vacuous} if
    $\exists I \in \mathcal I \setminus \mathcal G:\; T(I) = \mathsf{fail}$;
  \item \emph{satisfiable} if
    $\exists I \in \mathcal G:\; T(I) = \mathsf{pass}$;
  \item \emph{complete} if
    $\forall I \in \mathcal I:\; T(I) = \mathsf{pass}
      \iff I \in \mathcal G$.
\end{itemize}
\end{definition}

Thus:
\begin{description}
  \item[Non-vacuous] At least one incorrect implementation is actually
    detected:
    \[
      \exists I \in \mathcal I \setminus \mathcal G:\; T(I) = \mathsf{fail}.
    \]

  \item[Satisfiable] At least one intended implementation survives:
    \[
      \exists I \in \mathcal G:\; T(I) = \mathsf{pass}.
    \]

  \item[Complete] The suite perfectly separates intended from unintended
    implementations:
    \[
      \forall I \in \mathcal I:\;
      T(I) = \mathsf{pass} \iff I \in \mathcal G.
    \]
\end{description}

This is a \emph{counterfactual} view: a test suite classifies \emph{all}
implementations in $\mathcal I$, including those never actually deployed or
executed in any particular run. It asks which incorrect implementations
\emph{would} be rejected, not only what happens on the single implementation
actually in use.

In an architectural setting, a \emph{test component} together with its
surrounding harness induces such a test suite $T$ on $\mathcal I$: for each
implementation $I$, $T(I)$ records whether the composed system, with that
implementation plugged in, causes the test component to pass or fail.

Consequently, what we specify here is not ``a system that is already correct''
in isolation, but rather
\begin{quote}
a system whose interaction structure contains meaningful, adversarial tests
that separate intended implementations $\mathcal G$ from unintended ones
$\mathcal I \setminus \mathcal G$; correctness itself is then obtained only
under the additional assumption that the tests do in fact succeed on the
implementations actually deployed.
\end{quote}

\subsection{Why This Lens Is Useful}

The framework separates:
\begin{itemize}
\item specification semantics,
\item implementations,
\item joint behavior,
\item observation level,
\item architectural stance (solution concept),
\item global acceptability.
\end{itemize}

These parameters are orthogonal; small semantic changes can alter the architecture's stability landscape.
This fragility reveals mismatches between obligations, composition semantics, and global aims.

\subsection{Summary}

A tuple of component specifications induces a multi-player
game with binary demonic obligation semantics.
A solution concept $\mathsf{SC}$ selects profiles of architectural interest.
Comparing these with globally desirable behaviors yields precise notions of coupling, substitutability, specification minimality, and information hiding.

Some solution concepts (notably PNE) reveal undesired stability traps, while others may select desirable or optimal configurations.
The framework accommodates both without committing to a single interpretation.
